\documentclass[10pt,a4paper]{ctexart}
\usepackage[margin=2.0cm]{geometry}
\usepackage{microtype}
\usepackage{enumitem}

\setlength{\parindent}{2em}
\setlength{\parskip}{0.25em}
\setlist[itemize]{leftmargin=2.5em,itemsep=0.15em,topsep=0.2em,parsep=0pt}
\ctexset{section={beforeskip=1.1ex,afterskip=0.65ex}}

\title{课程项目选题报告\\[0.35em]\normalsize 面向有限计算资源的 RGB-D 关键帧选择与高效三维重建}
\author{姓名：黄远森\quad 学号：2023012174}
\date{2026 年 11 月 11 日}

\begin{document}
\pagestyle{plain}
\maketitle
\vspace{-3.0em}

\section{研究背景}
RGB-D 相机能够同步采集彩色图像和逐像素深度信息，是室内三维建模、机器人地图构建和增强现实的重要输入。常见重建流程会将连续深度帧融合到截断符号距离函数（TSDF）等体素表示中，但相邻帧通常存在大量冗余，全部处理会增加计算时间和数据传输开销。等间隔抽帧虽然简单，却可能在快速运动、视角变化、运动模糊或深度缺失时遗漏关键观测。因此，本项目拟研究一种适用于有限计算资源的轻量级关键帧选择方法，在减少输入帧数量的同时尽量保持三维重建质量。

\section{研究问题}
给定同步的 RGB-D 视频、相机内参和相机位姿，项目将从完整序列中选择少量关键帧用于三维融合。核心问题是：在相同帧数预算下，综合观测质量和三维几何信息的选择方法，能否比等间隔抽帧或单一运动阈值方法保留更高的重建质量？所提出的方法应当计算开销较低、无需训练模型，并能直接应用于公开 RGB-D 数据。项目将考察不同帧预算，以分析计算消耗与几何覆盖之间的权衡关系。

\section{拟采用的技术方案}
首先，对每个候选帧计算以下三类互补指标：
\begin{itemize}
  \item \textbf{观测质量：}利用拉普拉斯清晰度和有效深度比例，降低模糊帧及大面积深度缺失帧的优先级；
  \item \textbf{相机运动：}计算候选帧相对于近期关键帧的平移距离和旋转角度，识别有意义的视角变化；
  \item \textbf{几何新颖度：}将前一视角采样得到的三维点投影至候选视角，通过深度一致重叠率估计帧间冗余，并使用深度可靠性进行加权。
\end{itemize}

随后，将上述指标组合为候选帧信息得分，并在给定帧数预算下使用贪心方法完成选择。选择过程还将兼顾时间覆盖和相机位姿空间覆盖，避免关键帧集中在轨迹的局部区间。被选中的 RGB-D 帧将通过 Open3D 的可扩展 TSDF 方法进行融合，最终提取点云和三角网格。

实验计划采用公开 TUM RGB-D 数据集中的多个序列，并与全帧融合、等间隔抽帧和运动阈值抽帧进行比较。评测指标包括选帧比例、TSDF 融合时间、内存使用量、双向点云距离、重建完整度和几何 F-score。项目还将分别移除观测质量、相机运动和几何新颖度指标，以分析各组成部分的贡献。

\section{预期成果}
项目预期实现一个可复现的命令行原型：输入 TUM 格式 RGB-D 序列，输出关键帧、重建点云与网格、逐帧选择得分，以及重建质量与资源消耗的对比结果。预期通过实验得到清晰的质量与效率权衡关系，分析自适应关键帧选择在哪些场景和帧预算下能够稳定保留重要视角，并讨论不同策略在预算较充足时可能趋于一致的情况。最终成果包括完整源代码、配置文件、定量评测结果、课程报告和简短的三维重建演示视频。

\end{document}
